\begin{frame}
\frametitle{Type Syntax}

The (simplified) type syntax is defined as follows:
\begin{align*}
\tau &\doteq \ort{b}{\phi} ~|~ \urt{b}{\phi} ~|~ x{:}\tau \sarr \tau ~|~ x{:}\tau \wand \tau \\
\Gamma &\doteq \emptyset ~|~ x{:}\tau ~|~ \Gamma \sep \Gamma ~|~ \Gamma, \Gamma ~|~ \textcolor{gray}{\Code{emp}}
\end{align*}

\textbf{Comparison with SL:} 

\begin{itemize}
    \item $\ort{b}{\phi}$ is similar to an atomic predicate in pure logic.
    \item $\urt{b}{\phi}$ is similar to an atomic heap predicate in heap logic (e.g., $x \mapsto 4$).
    \item $x{:}\tau \sarr \tau$ is the standard (additive, nonlinear) arrow type.
    \item $x{:}\tau \wand \tau$ is similar to a magic wand (multiplicative, linear) in separation logic.
\end{itemize}

\textbf{Note:} For now, I am not sure if I should add $\textcolor{gray}{\Code{emp}}$ to the type context syntax.

\end{frame}

\begin{frame}
    \frametitle{Type Syntax}
    
    The basic bunched type (implication) only uses connectives to separate ``heap predicates'' and ``pure logic predicates.'' The key point is that there are two unit elements (empty contexts): $\emptyset$ and $\Code{emp}$, each for a different monoid.
    \begin{align*}
    x{:}\ort{b}{\phi} &\doteq x{:}\rt{b}{\phi}, \emptyset \\
    x{:}\urt{b}{\phi} &\doteq x{:}\rt{b}{\phi} \sep \Code{emp}
    \end{align*}

    \textbf{Examples:}
    \begin{align*}
    \forall a.\exists b.\forall c &\doteq a, (b \sep (c, \emptyset)) \\
    \forall a.\exists b.\exists c &\doteq a, (b \sep c) \\
    \exists a.\forall b.\exists c &\doteq a \sep (b, (c \sep \Code{emp})) \\
    \exists a.\forall b.\forall c &\doteq a \sep (b, c)
    \end{align*}
    
    \textbf{It is too complicated; thus, we still use a more ``SL-like'' notation.} 
    
\end{frame}

\begin{frame}
    \frametitle{Type Semantics}

\begin{align*}
    \denot{ \ort{b}{\phi} } &\doteq \{ e ~|~ \forall v. \mstep{e}{v} \impl \phi[\nu\mapsto v] \} \\
    \denot{ \urt{b}{\phi} } &\doteq \{ e ~|~ \forall v. \phi[\nu\mapsto v] \impl \mstep{e}{v}  \}
\end{align*}
    
\end{frame}

\begin{frame}
    \frametitle{Stuck Term and Divergence}

\textbf{Suresh \& Ben:} It is good to compare with Ramsey's work; it is also important to distinguish between a stuck term and divergence. One terminates in a non-value normal form, while the other does not terminate.

\textbf{Suresh:} An important ``reachability'' case to consider: $\lambda x.\; 1 \oplus (2 + \true)$. This function is not well-typed in the basic type system but is still reachable in some executions. (\textcolor{red}{ZZ: I don't believe this should be a typable program to consider, especially in ML-family languages.})

\textbf{Zhe:} This seems to conflict with the idea of ``refinement type,'' where if a term $e$ has refinement type $\tau$, it means $e$ also has $\eraserf{\tau}$ in the basic type system. We also need to change the standard reduction judgment $\mstep{e}{v}$; otherwise, we cannot distinguish between a stuck term and divergence.
\begin{align*}
    &\text{standard refinement type:} \forall v. \mstep{e}{v} \impl \phi[\nu\mapsto v] \\
    &\text{stuck term:} \forall N. \mstep{e}{N} \impl \neg \Code{isValue}(N) \\
    &\text{divergence term:} \forall N. \neg \mstep{e}{N}
\end{align*}

\end{frame}

\begin{frame}
    \frametitle{Stuck Term and Divergence}

\textbf{Try:} If we do not quite follow the manner of refinement types, what can we do?

\textbf{Zhe:} $\ort{b}{\bot}$, where $b$ is an arbitrary base type, can be used to represent an ``unreachable'' term, including both stuck terms and divergence, but it cannot distinguish between them.

\textbf{Ben:} However, the basic type system can type check the divergence term. Thus, if a term can have refinement type $\ort{b}{\bot}$ and basic type $b$ in the basic type system, it is a divergence term; otherwise, it is probably a stuck term (maybe Ramsey's work can be used to prove it must get stuck).

\textbf{Suresh:} Allowing divergence and tracking it is very annoying; it needs intersection types and dynamic choice depending on the calling context. \textcolor{red}{ZZ: I agree, in some sense divergence is an effect, it is context sensitive.}

\end{frame}

\begin{frame}
    \frametitle{Stuck Term and Divergence}

\textbf{Ben:} Can we use a two-judgment design (prove $\vdash$ and disprove $\nvdash$) to type check the stuck terms?

\textbf{Zhe:} It is not quite different from Ramsey's work.

\textbf{Ben:} Should we add three or more classes of base refinement types?

\textbf{Zhe:} It is kind of ugly and complex.

\textbf{Zhe:} In some sense, our paper is about ``incorrectness reasoning,'' separation, and resource. I believe it is quite far from Ramsey's work.

\end{frame}

\begin{frame}
    \frametitle{Stuck Term and Divergence}

End here.

\end{frame}

\begin{frame}
    \frametitle{Type Semantics}

\textbf{Consider a basic type $\Int\sarr\Int$. It can have the following refinement types:}

\begin{itemize}
    \item $x{:}\ort{\Int}{\vnu > 0}\sarr\ort{\Int}{\vnu > 0}$: For all positive inputs $x$, the function diverges or returns a positive result.
    \item $x{:}\ort{\Int}{\vnu > 0}\sarr\urt{\Int}{\vnu > 0}$: coverage type (similar to $\lambda x. \Code{alloc(posGen\;())}$ in SL).
    \item $x{:}\ort{\Int}{\vnu > 0}\wand\ort{\Int}{\vnu > 0}$: fallback to the normal arrow.
    \item $x{:}\ort{\Int}{\vnu > 0}\wand\urt{\Int}{\vnu > 0}$: fallback to the normal arrow.
\end{itemize}

\textbf{Compare with SL:} $x = 3 \wand y =  5$ fallback to $x = 3 \impl y = 5$.

\textbf{Note:} in Iris style triple, them are equal to:
\begin{align*}
    &\{x > 0 \land \textcolor{red}{\Code{emp}}\} y = ... \{ y > 0 \land \Code{emp} \} 
    \\&\{x > 0 \land \textcolor{red}{\Code{emp}}\} y = ... \{ y \mapsto \Code{posGen()} \} 
\end{align*}
Overapproximation clearly requires no resource.
    
\end{frame}

\begin{frame}
    \frametitle{Type Semantics}

\textbf{Consider a basic type $\Int\sarr\Int$. It can have the following refinement types:}

\begin{itemize}
    \item $x{:}\urt{\Int}{\vnu > 0}\wand\ort{\Int}{\vnu > 0}$: there exists a input such that the return is positive or diverges.
    \item $\lambda x. x - 100$ and $\Code{fix}\;f\;x. f\;x$ has this type; $\lambda x. 0 \oplus 1$ does not have this type.
    \item $x{:}\urt{\Int}{\vnu > 0}\wand\urt{\Int}{\vnu > 0}$: incorrectness triple (be careful with variable references).
    \item $\lambda x. x - 1$ has this type.
    \item $x{:}\urt{\Int}{\vnu > 0}\sarr\ort{\Int}{\vnu > 0}$: we require that the input must be able to cover all positive integers, and the result is positive (this application will not consume resources).
    \item $x{:}\urt{\Int}{\vnu > 0}\sarr\urt{\Int}{\vnu > 0}$: it is ill-defined, 
\end{itemize}

Requiring a ``resource'' but not consuming it is a little bit weird.
    
\end{frame}

\begin{frame}
    \frametitle{Type Semantics}

\textbf{To be concrete, what is the meaning of $x{:}\urt{b}{\phi}\sarr\ort{\Int}{\vnu > 0}$?}

\begin{itemize}
    \item $\lambda x. x - 100$ has type $x{:}\urt{\Int}{\top}\wand\ort{\Int}{\vnu > 0}$; if you give me a ``coverage guarantee,'' we now comsume it and return a positive result.
    \item $\lambda x. 1$ has type $x{:}\urt{\Int}{\top}\sarr\ort{\Int}{\vnu > 0}$; if you give me a ``coverage guarantee,'' I don't consume it and still return a positive result.
\end{itemize}

$\lambda x. 1$ can also have type $x{:}\ort{\Int}{\top}\sarr\ort{\Int}{\vnu > 0}$. In the CBV incorrectness reasoning setting, we have no way to check if a parameter has a ``coverage guarantee.''
    
\end{frame}

\begin{frame}
    \frametitle{Type Semantics: affine}

\textbf{Tricky case: $x{:}\urt{b}{\vnu > 0}\sarr\urt{\Int}{\vnu > 0}$?}

\textbf{SL: $\{x \mapsto 3\} e \{y \mapsto 3\}$?}

\begin{itemize}
    \item Standard SL: before execution I have a heap with singleton address $x$, after execution I have a heap with singleton address $y$. It may because $x = y$ and I do nothing, or $x \neq y$ and I free $x$ and allocate a new address $y$.
    \item Affine SL (e.g., Iris): I allow not to describe the state of address $x$ after execution, thus the $x$ can be not freed.
\end{itemize}

If we choose the first interpretation, it is the same with $x{:}\urt{b}{\vnu > 0}\wand\urt{\Int}{\vnu > 0}$.

If we choose the second interpretation, it is similar with $x{:}\ort{b}{\vnu > 0}\sarr\urt{\Int}{\vnu > 0}$.

In our pure functional setting, the function cannot modify the argument; the state of input argument is in charge of typing rules. We choose the interpretation similar with the second.
    
\end{frame}

\begin{frame}
    \frametitle{Type Semantics}

\textbf{It seems that we always have:}

\begin{itemize}
    \item $x{:}\ort{\Int}{\phi}\sarr\tau$: the normal arrow means over parameters.
    \item $x{:}\urt{\Int}{\phi}\wand\tau$: the magic wand means under parameters.
\end{itemize}
    
\end{frame} 

\begin{frame}[fragile]
    \frametitle{Resource Semantics: an SL-style Interpretation}

    \textbf{In order to explicitly show the meaning of resource in our setting:}

\textbf{A heap-based (modal logic style) interpretation:} We assume the under-type denotation is parameterized by a heap; coverage means there exists a heap makes it holds.

\textbf{Qualifier logic:} we allow $!x$ represents the value of address $x$ store in the heap.

\textbf{Type Context:} $x{:}\urt{b}{\phi}$ means some constraint about the assignment of $x$, and also \textbf{a singleton resource (heap)}.
\begin{align*}
    x{:}\urt{b}{\phi} \doteq \exists i{:}\Code{Addr}.\exists v{:}\Code{Choice}.\exists f{:}\Code{Choice}\sarr \rt{b}{\phi}. i\mapsto v \land x = f(v)
\end{align*}
Here, $i$ is an address, $v$ is a \emph{choice} that is stored at this address, and $f$ is a \textbf{surjective} function mapping a value of type $\Code{Choice}$ to (refined) type $b$ with qualifier $\phi$. 

Note that in order to make $f$ well-defined, we require $|\Code{Choice}| \geq |t|$ for any basic type $t$ (e.g., $|\Code{Choice}| = 2^{\aleph_0}$).

\end{frame}

\begin{frame}[fragile]
    \frametitle{Resource Semantics: an SL-style Interpretation}

\textbf{Type Context:} $x{:}\urt{b}{\phi}$ means some constraint about the assignment of $x$, and also \textbf{a singleton resource (heap)}.
\begin{align*}
    x{:}\urt{b}{\phi} \doteq \exists i{:}\Code{Addr}.\exists v{:}\Code{Choice}.\exists f{:}\Code{Choice}\sarr \rt{b}{\phi}. i\mapsto v \land x = f(v)
\end{align*}

\textbf{This means:} The value of $x$ depends on a heap fragment with address $i$; when the value $v$ is stored at this address, the value of $x$ is $f(v)$, since $f$ is surjective to the codomain $\rt{b}{\phi}$, the assignment of $x$ can cover $\phi$.

\textbf{Key:} The value of $x$ is fully determined by the heap fragment and the function $f$.

\end{frame}

\begin{frame}[fragile]
    \frametitle{Resource Semantics: an SL-style Interpretation}

\textbf{A heap-based interpretation:}
\begin{align*}
    x{:}\urt{b}{\phi}, y{:}\urt{b}{\phi} 
\end{align*}
means
\begin{align*}
    \exists v_1, v_2. (i_1 \mapsto v_1 \land x = f_1(v_1)) \land (i_2 \mapsto v_2 \land y = f_2(v_2))
\end{align*}
Here, the conjunction does not guarantee the separation of $i_1$ and $i_2$; or it assumes $f_1$ and $f_2$ share the same heap address, that is,
\begin{align*}
    \exists v. i \mapsto v \land x = f_1(v) \land y = f_2(v)
\end{align*}

\textbf{Key:} According to $\exists v. x = f_1(v)$, the assignment to $x$ can still cover the codomain of $f_1$, and similarly for $y$. However, $x$ and $y$ are not independent.
\end{frame}

\begin{frame}[fragile]
    \frametitle{Resource Semantics: a Formal Definition}

\textbf{A resource-based type denotation:} The denotation of a type context is a substitution \textbf{plus a set of heap addresses (or, abstractly, resource identifiers)}.
\begin{align*}
    \denot{\emptyset} &\doteq (\emptyset, \sigma_{\Code{id}}) \\
    \denot{x{:}\ort{b}{\phi}} &\doteq \{ (I, [x\mapsto v]) ~|~ v \in \denot{\ort{b}{\phi}} \} \\
    \denot{x{:}\urt{b}{\phi}} &\doteq \{ (\{i\}, [x\mapsto f(!i)]) ~|~ f \in \denot{\Code{Choice}\sarr \ort{b}{\phi}} \} \\
    \denot{\Gamma_1, \Gamma_2} &\doteq \{ (I, \sigma_1 \cup \sigma_2) ~|~ (I, \sigma_1) \in \denot{\Gamma_1} \land (I, \sigma_2) \in \denot{\textcolor{red}{\sigma_1(\Gamma_2)}} \} \\
    \denot{\Gamma_1 \sep \Gamma_2} &\doteq \{ (I_1 \cup I_2, \sigma_1 \cup \sigma_2) ~|~ 
    (I_1, \sigma_1) \in \denot{\Gamma_1} \land (I_2, \sigma_2) \in \denot{\Gamma_2} \land \textcolor{red}{I_1 \cap I_2 = \emptyset} \}
\end{align*}

\textbf{Footprints:}
\begin{itemize}
    \item The over type can bind to arbitrary resources since it is independent of any resources.
    \item The under type can only bind to a single resource identifier.
    \item The $,$ means the resources are exactly the same.
    \item The $\sep$ means the resources are disjoint; the denotation is only partially defined.
\end{itemize}
\end{frame}

\begin{frame}[fragile]
    \frametitle{Resource Semantics: a Formal Definition}

\textbf{Model:} $\Gamma \vdash e : \tau$ means for all $(I, \sigma) \in \denot{\Gamma}$, there exists a heap (resource) $h$ where $\Code{dom}(h) = I$. $h \models \vdash \sigma(e) : \sigma(\tau)$ holds.

\textbf{Examples:}
\begin{itemize}
    \item $\denot{x{:}\ort{\Int}{\vnu > 0}, y{:}\urt{\Int}{\vnu > 0}}$ has a singleton resource identifier (assume it is $i$)
    \begin{align*}
        &\{(\{i\}, [x\mapsto n, y\mapsto f(!i)]) ~|~ 
        n > 0 \land 
        \\&f \text{ is a surjective function covering all positive integers } \} 
    \end{align*}
    \item $\denot{x{:}\urt{\Int}{\vnu > 0}, y{:}\urt{\Int}{\vnu > 0}}$ also has a singleton resource identifier (assume it is $i$)
    \begin{align*}
        &\{(\{i\}, [x\mapsto f_1(!i), y\mapsto f_2(!i)]) ~|~ 
        f_1, f_2 \text{ are surjective functions covering all positive integers } \} 
    \end{align*}
    Note that the denotation considers all possible surjective functions $f_1$ and $f_2$; they can be equal. Thus, a judgment $e : \tau$ should hold even when $x = y$ (although the actual resource is unknown).
\end{itemize}
\end{frame}

\begin{frame}[fragile]
    \frametitle{Persistence}

\textbf{Question:} Can a coverage type $x{:}\ort{\Int}{\vnu > 0}\wand\urt{\Int}{\vnu > 0}$ be used (applied) multiple times?

\textbf{Answer:} It should be allowed, at least in the original paper.

\textbf{Question:} Can a resource be duplicated? Previously, we had no resource, and $f\;1$ receives one more resource (randomness).

\textbf{Answer:} Yes, and it should be duplicated. The function is a ``resource generator,'' and each time it is applied, it should provide a new resource.

\textbf{Problem:} This is not allowed by SL.

\textbf{Answer:} In Iris, we can add the persistence modality $\Box$ (similar to $!$ in linear logic).

\end{frame}

\begin{frame}[fragile]
    \frametitle{Persistence}

\textbf{Definition:} In Iris, $P$ is naturally persistent iff $P \vdash \Box P$ holds.

The hoare triple in Iris is persistent, which is the same with our function type.

\end{frame}

\begin{frame}[fragile]
    \frametitle{Typing Examples}

    \begin{minted}{ocaml}
        let p f g x =
            let y1 = f x in
            let y2 = g x in
            y2
    \end{minted}

Let
\begin{itemize}
    \item $\tau_f \doteq \urt{\Int}{\vnu > 0}\wand\urt{\Int}{\vnu > 0}$
    \item $\tau_g \doteq \ort{\Int}{\top}\sarr\urt{\Int}{\vnu > 0}$
\end{itemize}

We have:

\begin{prooftree}
    \hypo{
        f{:}\Box\tau_f, g{:}\Box\tau_g, x{:}\urt{\Int}{\top} \vdash \zlet{y1}{f\;x}{...} : \urt{\Int}{\vnu > 0}
    }
    \infer1[TFun]{
        \vdash p : f{:}\Box\tau_f \sarr g{:}\Box\tau_g \sarr x{:}\urt{\Int}{\top} \sarr \urt{\Int}{\vnu > 0} 
    }
\end{prooftree}
    
\end{frame}

\begin{frame}[fragile]
    \frametitle{Typing Examples}

Let
\begin{itemize}
    \item $\tau_f \doteq \urt{\Int}{\vnu > 0}\wand\urt{\Int}{\vnu > 0}$
    \item $\tau_g \doteq \ort{\Int}{\top}\sarr\urt{\Int}{\vnu > 0}$
\end{itemize}

We have:

\begin{prooftree}
    \hypo{
        x{:}\urt{\Int}{\top} \vdash x : \urt{\Int}{\top}
    }
    \infer1[TSub]{
        x{:}\urt{\Int}{\top} \vdash x : \urt{\Int}{\vnu > 0}
    }
\end{prooftree}

and,

\begin{prooftree}
    \hypo{
    }
    \infer1[TVar]{
        f{:}\Box\tau_f \vdash f : \tau_f
    }
\end{prooftree}

and,

\begin{prooftree}
    \hypo{
        (f{:}\Box\tau_f) \sep (g{:}\Box\tau_g, x{:}\urt{\Int}{\top}) \sep (g{:}\Box\tau_g, x{:}\ort{\Int}{\top}) \vdash ...
    }
    \infer1[TStruct]{
        f{:}\Box\tau_f, g{:}\Box\tau_g, x{:}\urt{\Int}{\top} \vdash ...
    }
\end{prooftree}
    
\end{frame}

\begin{frame}[fragile]
    \frametitle{Typing Examples}

Let
\begin{itemize}
    \item $\tau_f \doteq \urt{\Int}{\vnu > 0}\wand\urt{\Int}{\vnu > 0}$
    \item $\tau_g \doteq \ort{\Int}{\top}\sarr\urt{\Int}{\vnu > 0}$
\end{itemize}

We have:

\begin{prooftree}
    \hypo{
        g{:}\Box\tau_g, x{:}\ort{\Int}{\top} \vdash \zlet{y2}{g\;x}{y2} : \urt{\Int}{\vnu > 0}
    }
    \infer1[TStruct]{
        (g{:}\Box\tau_g, x{:}\ort{\Int}{\top}) \sep (y1{:}...) \vdash \zlet{y2}{g\;x}{y2} : \urt{\Int}{\vnu > 0}
    }
    \infer1[TApp]{
        (f{:}\Box\tau_f) \sep (g{:}\Box\tau_g, x{:}\urt{\Int}{\top}) \sep ... \vdash 
        \zlet{y1}{f\;x}{...} : \urt{\Int}{\vnu > 0}
    }
\end{prooftree}
    
\end{frame}

\begin{frame}[fragile]
    \frametitle{Higher-order Function Type}

Disprove a higher-order function will lead to ``under function parameter'' (higher-order symbolic execution).

    \begin{itemize}
        \item $x{:}\tau_f\sarr\tau$: for all function implementations with type $\tau_f$ ...
        \item $x{:}\ul{\tau_f}\wand\tau$: there exists a function implementation with type $\tau_f$ ...
    \end{itemize}

\textbf{What is the difference between ``forall'' and ``exists'' function parameters?}

\begin{minted}{ocaml}
let g (f: int -> int) = f 1
\end{minted}

\textbf{Over function parameter:}
\begin{itemize}
    \item $\vdash g : f{:}\Box(\urt{\Int}{\top} \wand \urt{\Int}{\top}) \sarr \urt{\Int}{\vnu = 3}$: for all possible functions having type $\urt{\Int}{\top} \wand \urt{\Int}{\top}$.
    \item We need to consider the function $\lambda x. x$, so $g$ cannot have this type.
\end{itemize}

\textbf{Under function parameter:}
\begin{itemize}
    \item $\vdash g : f{:}\Box(\urt{\Int}{\top} \wand \urt{\Int}{\top}) \wand \urt{\Int}{\vnu = 3}$: there exists a function implementation with type $\urt{\Int}{\top} \wand \urt{\Int}{\top}$.
    \item Only the function $\lambda x. x+2$ is considered, so $g$ can have this type.
\end{itemize}
    
\end{frame}

\begin{frame}[fragile]
    \frametitle{Persistence}
    
    \textbf{Higher-order symbolic execution:} The existential function parameter type is still persistent.
    
    \textbf{Example:} 
    \begin{itemize}
        \item In previous example,
        \begin{align*}
            \vdash g : (\urt{\Int}{\vnu > 0}\wand\urt{\Int}{\vnu > 0})\wand\urt{\Int}{\vnu = 3}
        \end{align*}
        \item Apply $f$ for multiple times doesn't change any things, ($y_1$ and $y_2$ are independent).
        \begin{minted}{ocaml}
        let g (f: int -> int) =
          let y1 = f (intGen ()) in
          let y2 = f (intGen ()) in
          y2
        \end{minted}
        \item If two application setting assuming two implementation $\lambda x.e_1$ and $\lambda x.e_2$, there always exists another function $\lambda x.e_1 \oplus e_2$.
        \item We loose this for base case, there is not a base value that can refine two distinct base values.
    \end{itemize}
    
    \end{frame}

\begin{frame}[fragile]
    \frametitle{Higher-order Function Type}

\textbf{Problem:} We use $\ort{b}{\phi}$ and $\urt{b}{\phi}$ to distinguished two modalities, how about the function types?

\begin{itemize}
    \item $x{:}\tau_f\sarr\tau$ is always overapproximated.
    \item $\ul{(x{:}\urt{\Int}{\top}\wand\urt{\Int}{\top})}$ means all possible implementations needs to be cover (a non-deterministic term).
    \item The $\Code{havoc}_{t}$ in higher-order symbolic execution.
\end{itemize}

\textbf{New syntax?}
\begin{align*}
    \tau \doteq \rt{b}{\phi} ~|~ \ul{\tau} ~|~ \ol{\tau} ~|~ \Box \tau ~|~ x{:}\sarr\tau ~|~ x{:}\wand\tau
\end{align*}

\textbf{One possible workaround:} Disallow under-style function types, and use forall function parameters instead.
\begin{itemize}
    \item $\vdash g : f{:}(\urt{\Int}{\top} \wand \urt{\Int}{\top} \interty \urt{\Int}{\vnu = 1} \wand \urt{\Int}{\vnu = 3}) \wand \urt{\Int}{\vnu = 3}$ (subtyping and covert to over-style function parameter).
\end{itemize}
    
\end{frame}

\begin{frame}[fragile]
    \frametitle{Higher-order Function Type}

\textbf{Problem:} We use $\ort{b}{\phi}$ and $\urt{b}{\phi}$ to distinguished two modalities, how about the function types?

\begin{itemize}
    \item $x{:}\tau_f\sarr\tau$ is always overapproximated.
    \item $\ul{(x{:}\urt{\Int}{\top}\wand\urt{\Int}{\top})}$ means all possible implementations needs to be cover (a non-deterministic term).
    \item The $\Code{havoc}_{t}$ in higher-order symbolic execution.
\end{itemize}

\textbf{New syntax?}
\begin{align*}
    \tau \doteq \rt{b}{\phi} ~|~ \ul{\tau} ~|~ \ol{\tau} ~|~ \Box \tau ~|~ x{:}\sarr\tau ~|~ x{:}\wand\tau
\end{align*}

\textbf{One possible workaround:} Disallow under-style function types, and use forall function parameters instead.
\begin{itemize}
    \item $\vdash g : f{:}(\urt{\Int}{\top} \wand \urt{\Int}{\top} \interty \urt{\Int}{\vnu = 1} \wand \urt{\Int}{\vnu = 3}) \wand \urt{\Int}{\vnu = 3}$ (subtyping and covert to over-style function parameter).
\end{itemize}
    
\end{frame}

% \begin{frame}[fragile]
%     \frametitle{Higher-order Function Type}

% However, it can help us distinguish between ``under function parameter'' (higher-order symbolic execution) and ``over function parameter'' (already checked functions).

%     \begin{itemize}
%         \item $x{:}\tau_f\sarr\tau$: for all function implementations with type $\tau_f$ ...
%         \item $x{:}\tau_f\wand\tau$: there exists a function implementation with type $\tau_f$ ...
%     \end{itemize}

% \textbf{What is the difference between ``forall'' and ``exists'' function parameters?}

% \begin{minted}{ocaml}
% let g (f: int -> int) = f 1
% \end{minted}

% \textbf{Over function parameter:}
% \begin{itemize}
%     \item $\vdash g : f{:}(\urt{\Int}{\vnu > 0} \wand \ort{\Int}{\vnu > 0}) \sarr \ort{\Int}{\vnu = 3}$: for all possible function implementations that have type $\urt{\Int}{\vnu > 0} \wand \ort{\Int}{\vnu > 0}$.
%     \item We need to consider the function $\lambda x. x$, so $g$ cannot have this type.
% \end{itemize}

% \textbf{Under function parameter:}
% \begin{itemize}
%     \item $\vdash g : f{:}(\urt{\Int}{\vnu > 0} \wand \ort{\Int}{\vnu > 0}) \wand \ort{\Int}{\vnu = 3}$: there exists a function implementation with type $\urt{\Int}{\vnu > 0} \wand \ort{\Int}{\vnu > 0}$.
%     \item Only the function $\lambda \_. 1$ is considered, so $g$ can have this type.
% \end{itemize}

% We will come back to this later.
    
% \end{frame}

% \begin{frame}[fragile]
%     \frametitle{Persistence}
    
%     \textbf{Question:} A concrete function implementation can only have a persistent type; can a function parameter be non-persistent?
    
%     \textbf{Base type:} A parameter $\urt{\Int}{\vnu > 0}$ is not persistent, since it corresponds to a random generator $\Code{posGen\;()}$ (instead of a value), so it will fall back to a concrete value after application.
    
%     \textbf{Higher-order symbolic execution:} The existential function parameter type is not persistent.
    
%     \textbf{Example:} 
%     \begin{itemize}
%         \item $f$ is a function such that ``there exists a function that can make its return be $3$'':
%         \begin{align*}
%             \vdash f : (\urt{\Int}{\vnu > 0}\wand\ort{\Int}{\vnu > 0})\wand\ort{\Int}{\vnu = 3}
%         \end{align*}
%         \item We cannot prove the following program can terminate:
%         \begin{minted}{ocaml}
%         let g (k1: int -> int) (k2: int -> int) = 
%           let _ = f k1 in
%           let _ = f k2 in
%           ()
%         \end{minted}
%     \end{itemize}
    
%     \end{frame}

% \begin{frame}[fragile]
%     \frametitle{Reconsidering Overapproximation}
    
%     \textbf{Intuition:} It seems that ``persistency'' is deeply connected with overapproximation.
    
%     \begin{align*}
%         \denot{\emptyset} &\doteq (\emptyset, \sigma_{\Code{id}}) \\
%         \denot{x{:}\rt{b}{\phi}} &\doteq \{ (\{i\}, [x\mapsto f(!i)]) ~|~ \exists h, h \models \phi[\vnu \mapsto f(!i)] \} \\
%         \denot{x{:}\urt{b}{\phi}} &\doteq \{ (\{i\}, [x\mapsto f(!i)]) ~|~ \forall v. \phi[\vnu \mapsto v] \impl \exists h, h\models f(!i) = v \} \\
%         \denot{x{:}\ort{b}{\phi}} &\doteq \denot{x{:}\Box \rt{b}{\phi}} \\
%         \denot{\Gamma_1, \Gamma_2} &\doteq \{ (I, \sigma_1 \cup \sigma_2) ~|~ (I, \sigma_1) \in \denot{\Gamma_1} \land (I, \sigma_2) \in \denot{\textcolor{red}{\sigma_1(\Gamma_2)}} \} \\
%         \denot{\Gamma_1 \sep \Gamma_2} &\doteq \{ (I_1 \cup I_2, \sigma_1 \cup \sigma_2) ~|~ 
%             (I_1, \sigma_1) \in \denot{\Gamma_1} \land (I_2, \sigma_2) \in \denot{\Gamma_2} \land \textcolor{red}{I_1 \cap I_2 = \emptyset} \}
%     \end{align*}
    
% \end{frame}

% \begin{frame}[fragile]
%     \frametitle{Problem: Function Parameters}
    
%     \begin{itemize}
%         \item $f{:}\Box\tau_f\sarr \tau$: overapproximated function parameter.
%         \item $f{:}\tau_f\wand \tau$: underapproximated function parameter.
%     \end{itemize}

%     \textbf{Flip operation:} Flip a function type?
% \begin{align*}    
%     x{:}\urt{b}{\phi} \sep x{:}\ort{b}{\phi} \equiv x{:}\urt{b}{\phi}
% \end{align*}

% \textbf{Negate operation:} Negation also seems to remove the persistency modality.
    
% \end{frame}


\begin{frame}[fragile]
    \frametitle{Problem: divergence Reasoning}
    
    \begin{itemize}
        \item $\ort{\Int}{\bot}$ means divergence. 
        \item Every function has a default type $\ort{\Int}{\bot} \sarr \ort{\Int}{\bot}$, stronger than $\ort{\Int}{\bot} \sarr \ort{\Int}{\top}$.
    \end{itemize}

    \textbf{Example:} $\vdash f : \ort{\Int}{\vnu > 0} \sarr \ort{\Int}{\bot}$.
    \begin{minted}{ocaml}
        let rec f (x: int) = f x

        let rec f (x: int) = 
          if x <= 0 then 0 else f (x + 1)
    \end{minted}
    \begin{align*}
        &x{:}\ort{\Int}{\vnu > 0}, \_{:}\ort{\Unit}{x \leq 0} \vdash 0 : \urt{\Int}{\bot} 
        \\&\quad \text{ since type context has no inheritant. }
    \end{align*}
    
\end{frame}